\documentclass[10pt,a4paper]{amsart}
\usepackage[margin=19mm]{geometry}
\usepackage{amsmath,amssymb,mathtools}
\usepackage[scr=rsfs,cal=euler]{mathalfa}
\usepackage{xfrac}
\usepackage[unicode,hidelinks,bookmarks=false]{hyperref}
\newcommand{\ie}{i.e.\ }
\newcommand{\cf}{cf.\ }
\newcommand{\resp}{respectively\ }
\newcommand{\Subsection}{\subsection}
\providecommand{\nobreakdash}{\nobreak\hbox{-}}
\setlength{\emergencystretch}{2em}
\multlinegap0pt
\usepackage{bm}
\usepackage{bbm}
\usepackage{dsfont}
\usepackage{xspace}
\usepackage{cancel}
\usepackage{mathtools}
\usepackage{amsthm}
\makeatletter
\@ifundefined{theorem}{\newtheorem{theorem}{Theorem}}{}
\@ifundefined{definition}{\newtheorem{definition}[theorem]{Definition}}{}
\@ifundefined{lemma}{\newtheorem{lemma}[theorem]{Lemma}}{}
\@ifundefined{remark}{\newtheorem{remark}[theorem]{Remark}}{}
\makeatother
\title[Forest Polynomials and Pattern Avoidance]{Forest Polynomials and Pattern Avoidance}
\author{Annie Guo, Dora Woodruff}
\date{Source: arXiv:2602.04036v1 (CC BY 4.0)}
\begin{document}
\maketitle

\noindent\emph{Source and licence.} Forest Polynomials and Pattern Avoidance. Version arXiv:2602.04036v1 (CC BY 4.0), distributed under the Creative Commons Attribution 4.0 International licence, as recorded in the arXiv licence metadata for this exact version and shown on the abstract page https://arxiv.org/abs/2602.04036v1. Full rights notice: licenses/arxiv-2602.04036-CC-BY-4.0.txt.\par\smallskip

\noindent\textit{Editorial scope.} Counted unit: the Lemma whose source label is \texttt{lemma:sufficientcase2} in the article (source0.tex lines 1068--1070, source label \texttt{lemma:sufficientcase2}), delivered together with the article's own complete original proof of it (source0.tex lines 1071--1103, one \texttt{proof} environment of 33 lines). No mathematics is written, repaired or re-derived: statement and proof are the article's own text line for line. Required context retained uncounted: definition:matching (lines 656--662); rmk:diagonalremark (lines 739--741). Every definition, display and label of those lines is retained unchanged. Omitted from this package and not redistributed: 1--655, 663--738, 742--1067, 1104--1114. Nothing inside a retained excerpt is omitted or altered: every retained body is delivered exactly as the article's lines are written, so no omission receipt of any kind accompanies this package. Citations: the retained excerpts carry no citation command at all, so no bibliography entry of the article is transcribed and no reference is redistributed. Reference adaptation: none was needed and none was made; every reference inside the delivered unit resolves inside it, so no unresolved reference or citation is produced. Printed slips kept literal: no printed word, symbol, hypothesis or number of the counted statement or of its proof was altered. Layout and class: the article's own class is \texttt{article} with packages including graphicx, mathrsfs, tikz-cd, hyperref, amsmath, quiver, amsthm, amssymb, comment, parskip, geometry, ytableau, multicol, tikz; the wrapper is the lane's pinned \texttt{amsart} editorial wrapper at 10pt on A4, which restates the theorem-like environments the retained text needs and adds the article's own macro definitions that the retained text needs (none), with extra wrapper packages (bm, bbm, dsfont, xspace, cancel, mathtools). The delivered numbering is the unit's own. Assets: no retained span contains a figure environment, a graphics-inclusion command, a PDF-page inclusion, a table or a TikZ picture; no image file is copied into this package, the package declares no asset dependency and no image file is redistributed with it. Licence: version arXiv:2602.04036v1 of this article is distributed under the Creative Commons Attribution 4.0 International licence, as recorded in the arXiv licence metadata for this exact version and shown on the abstract page https://arxiv.org/abs/2602.04036v1. A full rights notice is delivered at licenses/arxiv-2602.04036-CC-BY-4.0.txt. Characters: every retained excerpt is pure ASCII, so the delivered engine, XeLaTeX with its default Latin Modern text font, reports no missing character.
\begin{definition}[Covering crossings]\label{definition:matching}
   We recursively define a certain relation between crossings of $P_{\text{bott}}(w)$ as follows. First, consider row $1$. Find the first nonempty row $j$ below row $1$; then, the $(j-1)$st crossing of row $1$ \textit{covers} the rightmost crossing of row $j$. (The first $(j-2)$ crossings of row $1$ do not cover anything). We say that the first $j-1$ crossings of row $1$ are all `completed.'  

   Recursively, apply the same procedure to row $j$. Only return to row $1$ when all crossings in row $j$ have been `completed.' 
   
    Then, when all crossings in row $1$ have been completed, find the next nonempty row $i$ with non-completed crossings, and start from the beginning, with $i$ playing the same role as $1$. If there is no such row $i$, we are done. 
\end{definition}

\begin{remark}\label{rmk:diagonalremark}
Suppose we have crossings $c_w,c_v$ corresponding to vertices $w,v$. Then if $c_w$ lies in the diagonal $d_v-1$ or greater in $P_{\text{bott}}(w)$ (equivalently $c_v$ lies in the diagonal $d_w+1$ or smaller), $v$ is in the subtree of $w$. 
\end{remark}

\begin{lemma}\label{lemma:sufficientcase2}
    Suppose that $(c_w,c_v)$ is a bad parent-child pair. Then, suppose the rightmost crossing of row $r_v$ can slide up past $r_w$ only after some box in row $r_w$ makes a simple ladder move. (That is, there exists a crossing on diagonal $d_v$ or $d_{v}-1$). Then, $w$ contains either a $1432$, $2431$ or a $14523$ pattern.
\end{lemma}

\begin{proof}
    Suppose that $c_{v'}$ is our crossing on diagonal $d_v$ or $d_{v-1}$, between rows $r_v$ and $r_w$. First, suppose that $c_{v'}$ is not \textit{in} row $r_w$ (in $P_{\text{bott}}(w)$). But then, by Remark \ref{rmk:diagonalremark}, $c_w$ cannot directly be a parent of $c_v$, contradicting the assumptions of our lemma. (And inductively, we could keep searching for a bad parent-child pair by letting $c_{v'}$ play the role of $c_v$ instead).  
    

    Therefore, we may assume that the only crossing on diagonal $d_v$ or $d_{v}-1$ between rows $r_v$ and $r_w$ is contained in row $r_w$ itself. In this case, notice that there must exist some row $r_1$ above $r_w$ satisfying \[w(r_1)<w(r_w),w(r_v)\] 
    (If not, then no crossing in $r_w$ can ever move starting from $P_{\text{bott}}(w)$, and $c_v$ can never slide into row $r_w$). This $r_1$ will play the role of the $1$ in our desired patterns. 
    
    Now, we break into two subcases:

    \textbf{Subcase 1:} There is a crossing in $r_w$ on diagonal $d_v$ (as well as $d_{v}-1$, since $P_{\text{bott}}(w)$ is left-adjusted). 

    In this case, $w$ contains a $1432$ pattern or a $2431$ pattern, formed by indices $r_1, r_w, r_v$, and some row $x > r_v$. More specifically, there must be an $x > r_v$ with $w(x) < w(r_v)$, so that $L(r_v) \neq 0$. Furthermore, if $w(r_1) < (x)$, we have a $1432$ pattern. If $r_1 > x$, we have a $2431$ pattern. 
    
    Finally, consider:

    \textbf{Subcase 2:} row $r_w$ contains a crossing in $d_v-1$, but not in $d_v$. 

    In this case, we claim that $w$ contains a $14523$ pattern or a $1432$ pattern. Because $c_v$ can slide along its diagonal until it is directly beneath the rightmost crossing of row $r_w$, we have $w(r_w) < w(r_v)$.

    Now, we claim that row $r_v$ must contain more than one crossing. If it only contains $c_v$, then Lemma \ref{definition:matching} implies that $c_v$ is actually the right child of the rightmost crossing in row $r_w$. But this is impossible, since we said the diagonals of $c_v$ and $c_w$ must be more than $1$ apart. 

    Therefore, $L(r_v) \geq 2$, so there are two more rows $x, y$ with $x, y > r_v$ and $w(r_v) > w(x), w(y)$. Finally, depending on the six possible relative orderings of $w(r_1), w(x),$ and $w(y)$, indices $r_1, r_w, r_v, x, y$ give us one of the following patterns in $w$: 
    \begin{enumerate}
        \item $15432$
        \item $25413$
        \item $25431$
        \item $35412$
        \item $35421$ or 
        \item $15423$
    \end{enumerate}

    In the first three cases, there is a $1432$ subpattern. In the next two cases, there is a $2431$ subpattern. And in the last case, we have our last forbidden pattern, $15423$. 
\end{proof}

\end{document}
